\pagestyle{empty}
\tight
\begin{tblr}{width=\textwidth,colspec={|Q[c,m,1.9cm]|X[l,m]|},hlines,vlines,hspan=minimal,row{1}={bg=headblue},rowsep=1pt,colsep=3pt}
\SetCell{c,m}\hfil\logo\hfil & \textbf{POLITEKNIK NEGERI MALANG}\par\textbf{JURUSAN TEKNOLOGI INFORMASI}\par\textbf{PROGRAM STUDI: D4 TEKNIK INFORMATIKA} \\
\SetCell[c=2]{c} \textbf{RENCANA PEMBELAJARAN SEMESTER (RPS)} & \\
\end{tblr}
\begin{longtblr}{width=\textwidth,colspec={|Q[l,m,1.9cm]|X[0.85,l,m]|X[1.45,l,m]|X[0.75,l,m]|X[1.15,l,m]|},hlines,vlines,hspan=minimal,rowsep=1pt,colsep=2pt,row{1,3}={bg=headgray,font=\bfseries}}
MATA KULIAH & KODE & BOBOT (SKS)/jam & SEMESTER & TGL. PENYUSUNAN \\
Business Intelligence & RTI255005 & 2 SKS/4 jam & 5 & 1 Juli 2025 \\
OTORISASI & Dosen Pengembang RPS & Koordinator RMK & \SetCell[c=2]{l} Ka PRODI &  \\
 & Farid Angga Pribadi, S.Kom., M.Kom. & Farid Angga Pribadi, S.Kom., M.Kom. & \SetCell[c=2]{l}{\vspace{8mm}\par Dr. Ely Setyo Astuti, S.T., M.T.} &  \\
\SetCell[r=11]{l} Capaian Pembelajaran (CP) & \SetCell[c=4]{l,bg=headgray,font=\bfseries} Capaian Pembelajaran Lulusan Program Studi (CPL-Prodi) &  &  &  \\
 & CPL07 & \SetCell[c=3]{l} Mampu mengembangkan sistem komputasi cerdas dan melakukan analisis data berbasis kecerdasan buatan untuk pengambilan keputusan. &  &  \\
 & CPL10 & \SetCell[c=3]{l} Mampu menganalisis persoalan kompleks dengan wawasan multidisiplin untuk menghasilkan solusi di bidang informatika dan ilmu komputer. &  &  \\
 & \SetCell[c=4]{l,bg=headgray,font=\bfseries} Capaian Pembelajaran mata kuliah (CPL-MK) &  &  &  \\
 & CPMK0705 & \SetCell[c=3]{l} Mahasiswa mampu melakukan pengolahan data dan visualisasi untuk menghasilkan insight yang bermakna bagi pengambilan keputusan bisnis. &  &  \\
 & CPMK0706 & \SetCell[c=3]{l} Mahasiswa mampu mengembangkan sistem BI berbasis AI/ML untuk analisis data dan pendukung keputusan. &  &  \\
 & CPMK1004 & \SetCell[c=3]{l} Mahasiswa mampu merancang dan menginterpretasikan visualisasi grafis data untuk analisis dan keputusan strategis. &  &  \\
 & \SetCell[c=4]{l,bg=headgray,font=\bfseries} Kemampuan akhir tiap tahapan belajar (Sub-CPMK) &  &  &  \\
 & SCPMK0705-04101 & \SetCell[c=3]{l} Mahasiswa mampu merancang dan mengimplementasikan data warehouse dengan skema bintang/snowflake dan proses ETL. &  &  \\
 & SCPMK0706-04102 & \SetCell[c=3]{l} Mahasiswa mampu membangun dashboard BI interaktif yang mengintegrasikan analisis OLAP dan visualisasi data. &  &  \\
 & SCPMK1004-04103 & \SetCell[c=3]{l} Mahasiswa mampu mengevaluasi kualitas visualisasi data dan menyajikan insight bisnis yang akurat kepada pemangku kepentingan. &  &  \\
\SetCell[r=2]{l} Penilaian & \SetCell[c=4]{l} *Nilai CPMK didapat dari agregasi nilai SUB-CPMK yang dibebankan &  &  &  \\
 & \SetCell[c=4]{l}{\tiny\begin{tblr}{width=\linewidth,colspec={|Q[c,m,0.1\linewidth]|Q[c,m,0.16\linewidth]|X[l,m]|X[l,m]|X[l,m]|X[l,m]|X[l,m]|X[l,m]|Q[c,m,0.08\linewidth]|},hlines,vlines,hspan=minimal,rowsep=1pt,colsep=0.7pt,row{1,2}={bg=headgray,font=\bfseries}}
Kode CPMK & Kode SCPMK & \SetCell[c=6]{c} Bobot per Bentuk Penilaian & & & & & & Total Bobot \\
 & & Tugas & Kuis 1 & UTS & Kuis 2 & PBL & UAS & \\
CPMK0705 & SCPMK0705-04101 & 12\%: data warehouse dan ETL & 5\%: data warehouse dan ETL & 7\%: DW, ETL, dan OLAP & 0\% & 0\% & 0\% & 24\% \\
CPMK0706 & SCPMK0706-04102 & 10\%: tools BI dan dashboard & 0\% & 6\%: DW, ETL, dan OLAP & 0\% & 10\%: implementasi dashboard & 12\%: presentasi dashboard & 38\% \\
CPMK1004 & SCPMK1004-04103 & 8\%: visualisasi data & 0\% & 0\% & 5\%: visualisasi dan dashboard & 12\%: implementasi dashboard & 13\%: presentasi dashboard & 38\% \\
\SetCell[c=8]{r}\textbf{Total Per Penilaian} & & & & & & & & \textbf{100\%} \\
\end{tblr}} &  &  &  \\
Diskripsi Singkat Mata Kuliah & \SetCell[c=4]{l} Mata kuliah Business Intelligence membangun kemampuan merancang data warehouse, proses ETL, analisis OLAP, dan visualisasi dashboard interaktif untuk menghasilkan insight bisnis yang mendukung pengambilan keputusan strategis. &  &  &  \\
Bahan Kajian / Materi Pembelajaran & \SetCell[c=4]{l}\begin{enumerate}\item Pengantar BI, data warehouse, dan proses ETL\item Analisis OLAP dan operasi multidimensional\item Visualisasi data, tools BI, dan desain dashboard\item Data mining, predictive analytics, dan real-time BI\item Implementasi proyek dashboard BI dan evaluasi insight\end{enumerate} &  &  &  \\
\SetCell[r=2]{l} Pustaka & \SetCell[c=4]{l} \textbf{Utama:}\begin{enumerate}\item Kimball, R. \& Ross, M. 2013. The Data Warehouse Toolkit (3rd ed.). Wiley.\item Few, S. 2012. Show Me the Numbers. Analytics Press.\item Dokumentasi Power BI dan Tableau resmi.\end{enumerate} &  &  &  \\
 & \SetCell[c=4]{l} \textbf{Pendukung:} Materi LMS, dokumentasi resmi perangkat lunak, dan jobsheet praktikum. &  &  &  \\
Media Pembelajaran & \SetCell[c=2]{l} \textbf{Software:} Power BI, Tableau, Metabase, Python, LMS. &  & \SetCell[c=2]{l} \textbf{Hardware:} Komputer/laptop, proyektor. &  \\
Nama Dosen Pengampu & \SetCell[c=4]{l} Tim Pengajar Program Studi D4 TI &  &  &  \\
Matakuliah Syarat & \SetCell[c=4]{l} - &  &  &  \\
\end{longtblr}
\begin{longtblr}{width=\textwidth,colspec={|Q[c,m,0.06\textwidth]|X[1.35,l,m]|X[1.08,l,m]|X[1.16,l,m]|X[0.7,l,m]|X[1.24,l,m]|X[1.06,l,m]|X[0.98,l,m]|Q[c,m,0.05\textwidth]|},hlines,vlines,hspan=minimal,rowhead=2,row{1,2}={bg=headgreen,font=\bfseries},rowsep=1pt,colsep=1.5pt}
Mgg.\par Ke & Kemampuan Akhir Yang Direncanakan (Sub-CP-MK) & Bahan kajian (Materi Pembelajaran) & Bentuk dan Metode Pembelajaran & Estimasi Waktu & Pengalaman Belajar Mahasiswa & Kriteria \& Bentuk Penilaian & Indikator Penilaian & Bobot Penilaian (\%) \\
(1) & (2) & (3) & (4) & (5) & (6) & (7) & (8) & (9) \\
1 & SCPMK0705-04101 & Pengantar BI: konsep, arsitektur, dan ekosistem. & Modalitas: Blended Learning. Bentuk: Luring/Daring. Metode: Case Method, studi kasus BI, analisis data, dan presentasi kelompok. & 1 x 2 x 50' tatap muka; 1 x 2 x 50' tugas/praktik mandiri. & Mahasiswa mengerjakan aktivitas terarah untuk pengantar BI dan menunjukkan evidence hasil belajar. & Bentuk penilaian: Pengantar BI. Mengacu rubrik RTM. & Ketepatan konsep; kualitas analisis data; validasi dan dokumentasi. & 2 \\
2 & SCPMK0705-04101 & Data warehouse: skema bintang dan dimensi-fakta. & Modalitas: Blended Learning. Bentuk: Luring/Daring. Metode: Case Method, studi kasus BI, analisis data, dan presentasi kelompok. & 1 x 2 x 50' tatap muka; 1 x 2 x 50' tugas/praktik mandiri. & Mahasiswa mengerjakan aktivitas terarah untuk data warehouse dan menunjukkan evidence hasil belajar. & Bentuk penilaian: Data Warehouse. Mengacu rubrik RTM. & Ketepatan konsep; kualitas analisis data; validasi dan dokumentasi. & 3 \\
3 & SCPMK0705-04101 & ETL: Extract, Transform, Load. & Modalitas: Blended Learning. Bentuk: Luring/Daring. Metode: Case Method, studi kasus BI, analisis data, dan presentasi kelompok. & 1 x 2 x 50' tatap muka; 1 x 2 x 50' tugas/praktik mandiri. & Mahasiswa mengerjakan aktivitas terarah untuk ETL dan menunjukkan evidence hasil belajar. & Bentuk penilaian: ETL. Mengacu rubrik RTM. & Ketepatan konsep; kualitas analisis data; validasi dan dokumentasi. & 3 \\
4 & SCPMK0705-04101 & OLAP: slice, dice, drill-down, dan roll-up. & Modalitas: Blended Learning. Bentuk: Luring/Daring. Metode: Case Method, studi kasus BI, analisis data, dan presentasi kelompok. & 1 x 2 x 50' tatap muka; 1 x 2 x 50' tugas/praktik mandiri. & Mahasiswa mengerjakan aktivitas terarah untuk OLAP dan menunjukkan evidence hasil belajar. & Bentuk penilaian: OLAP. Mengacu rubrik RTM. & Ketepatan konsep; kualitas analisis data; validasi dan dokumentasi. & 3 \\
5 & SCPMK0705-04101 & Kuis 1: data warehouse dan ETL. & Modalitas: Blended Learning. Bentuk: Luring/Daring. Metode: Case Method, studi kasus BI, analisis data, dan presentasi kelompok. & 1 x 2 x 50' tatap muka; 1 x 2 x 50' tugas/praktik mandiri. & Mahasiswa mengerjakan aktivitas terarah untuk kuis 1 dan menunjukkan evidence hasil belajar. & Bentuk penilaian: Kuis 1. Mengacu rubrik RTM. & Ketepatan konsep; kualitas analisis data; validasi dan dokumentasi. & 5 \\
6 & SCPMK1004-04103 & Visualisasi data: prinsip dan jenis chart. & Modalitas: Blended Learning. Bentuk: Luring/Daring. Metode: Case Method, studi kasus BI, analisis data, dan presentasi kelompok. & 1 x 2 x 50' tatap muka; 1 x 2 x 50' tugas/praktik mandiri. & Mahasiswa mengerjakan aktivitas terarah untuk visualisasi data dan menunjukkan evidence hasil belajar. & Bentuk penilaian: Visualisasi Data. Mengacu rubrik RTM. & Ketepatan konsep; kualitas analisis data; validasi dan dokumentasi. & 3 \\
7 & SCPMK0706-04102 & Tools BI: Power BI, Tableau, dan Metabase. & Modalitas: Blended Learning. Bentuk: Luring/Daring. Metode: Case Method, studi kasus BI, analisis data, dan presentasi kelompok. & 1 x 2 x 50' tatap muka; 1 x 2 x 50' tugas/praktik mandiri. & Mahasiswa mengerjakan aktivitas terarah untuk tools BI dan menunjukkan evidence hasil belajar. & Bentuk penilaian: Tools BI. Mengacu rubrik RTM. & Ketepatan konsep; kualitas analisis data; validasi dan dokumentasi. & 3 \\
8 & SCPMK0705-04101, SCPMK0706-04102 & UTS: data warehouse, ETL, dan OLAP. & Modalitas: Blended Learning. Bentuk: Luring/Daring. Metode: Case Method, studi kasus BI, analisis data, dan presentasi kelompok. & 1 x 2 x 50' tatap muka; 1 x 2 x 50' tugas/praktik mandiri. & Mahasiswa mengerjakan aktivitas terarah untuk UTS dan menunjukkan evidence hasil belajar. & Bentuk penilaian: UTS. Mengacu rubrik RTM. & Ketepatan konsep; kualitas analisis data; validasi dan dokumentasi. & 13 \\
9 & SCPMK1004-04103 & Dashboard design: KPI dan storytelling data. & Modalitas: Blended Learning. Bentuk: Luring/Daring. Metode: Case Method, studi kasus BI, analisis data, dan presentasi kelompok. & 1 x 2 x 50' tatap muka; 1 x 2 x 50' tugas/praktik mandiri. & Mahasiswa mengerjakan aktivitas terarah untuk dashboard design dan menunjukkan evidence hasil belajar. & Bentuk penilaian: Dashboard Design. Mengacu rubrik RTM. & Ketepatan konsep; kualitas analisis data; validasi dan dokumentasi. & 3 \\
10 & SCPMK0706-04102 & Data mining: klasifikasi dan asosiasi. & Modalitas: Blended Learning. Bentuk: Luring/Daring. Metode: Case Method, studi kasus BI, analisis data, dan presentasi kelompok. & 1 x 2 x 50' tatap muka; 1 x 2 x 50' tugas/praktik mandiri. & Mahasiswa mengerjakan aktivitas terarah untuk data mining dan menunjukkan evidence hasil belajar. & Bentuk penilaian: Data Mining. Mengacu rubrik RTM. & Ketepatan konsep; kualitas analisis data; validasi dan dokumentasi. & 3 \\
11 & SCPMK0706-04102 & Predictive analytics dalam BI. & Modalitas: Blended Learning. Bentuk: Luring/Daring. Metode: Case Method, studi kasus BI, analisis data, dan presentasi kelompok. & 1 x 2 x 50' tatap muka; 1 x 2 x 50' tugas/praktik mandiri. & Mahasiswa mengerjakan aktivitas terarah untuk predictive analytics dan menunjukkan evidence hasil belajar. & Bentuk penilaian: Predictive Analytics. Mengacu rubrik RTM. & Ketepatan konsep; kualitas analisis data; validasi dan dokumentasi. & 3 \\
12 & SCPMK1004-04103 & Kuis 2: visualisasi dan dashboard BI. & Modalitas: Blended Learning. Bentuk: Luring/Daring. Metode: Case Method, studi kasus BI, analisis data, dan presentasi kelompok. & 1 x 2 x 50' tatap muka; 1 x 2 x 50' tugas/praktik mandiri. & Mahasiswa mengerjakan aktivitas terarah untuk kuis 2 dan menunjukkan evidence hasil belajar. & Bentuk penilaian: Kuis 2. Mengacu rubrik RTM. & Ketepatan konsep; kualitas analisis data; validasi dan dokumentasi. & 5 \\
13 & SCPMK0706-04102 & Real-time BI dan streaming analytics. & Modalitas: Blended Learning. Bentuk: Luring/Daring. Metode: Case Method, studi kasus BI, analisis data, dan presentasi kelompok. & 1 x 2 x 50' tatap muka; 1 x 2 x 50' tugas/praktik mandiri. & Mahasiswa mengerjakan aktivitas terarah untuk real-time BI dan menunjukkan evidence hasil belajar. & Bentuk penilaian: Real-time BI. Mengacu rubrik RTM. & Ketepatan konsep; kualitas analisis data; validasi dan dokumentasi. & 3 \\
14 & SCPMK0705-04101 & Self-service BI dan data governance. & Modalitas: Blended Learning. Bentuk: Luring/Daring. Metode: Case Method, studi kasus BI, analisis data, dan presentasi kelompok. & 1 x 2 x 50' tatap muka; 1 x 2 x 50' tugas/praktik mandiri. & Mahasiswa mengerjakan aktivitas terarah untuk self-service BI dan menunjukkan evidence hasil belajar. & Bentuk penilaian: Self-service BI. Mengacu rubrik RTM. & Ketepatan konsep; kualitas analisis data; validasi dan dokumentasi. & 3 \\
15 & SCPMK0706-04102, SCPMK1004-04103 & Studi kasus BI industri. & Modalitas: Blended Learning. Bentuk: Luring/Daring. Metode: Case Method, studi kasus BI, analisis data, dan presentasi kelompok. & 1 x 2 x 50' tatap muka; 1 x 2 x 50' tugas/praktik mandiri. & Mahasiswa mengerjakan aktivitas terarah untuk studi kasus BI dan menunjukkan evidence hasil belajar. & Bentuk penilaian: Studi Kasus BI. Mengacu rubrik RTM. & Ketepatan konsep; kualitas analisis data; validasi dan dokumentasi. & 3 \\
16 & SCPMK0706-04102, SCPMK1004-04103 & PBL: implementasi dashboard BI. & Modalitas: Blended Learning. Bentuk: Luring/Daring. Metode: Case Method, studi kasus BI, analisis data, dan presentasi kelompok. & 1 x 2 x 50' tatap muka; 1 x 2 x 50' tugas/praktik mandiri. & Mahasiswa mengerjakan aktivitas terarah untuk PBL dashboard BI dan menunjukkan evidence hasil belajar. & Bentuk penilaian: PBL Dashboard BI. Mengacu rubrik RTM. & Ketepatan konsep; kualitas analisis data; validasi dan dokumentasi. & 22 \\
17 & SCPMK0706-04102, SCPMK1004-04103 & UAS: presentasi dan demo dashboard BI. & Modalitas: Blended Learning. Bentuk: Luring/Daring. Metode: Case Method, studi kasus BI, analisis data, dan presentasi kelompok. & 1 x 2 x 50' tatap muka; 1 x 2 x 50' tugas/praktik mandiri. & Mahasiswa mengerjakan aktivitas terarah untuk UAS dan menunjukkan evidence hasil belajar. & Bentuk penilaian: UAS. Mengacu rubrik RTM. & Ketepatan konsep; kualitas analisis data; validasi dan dokumentasi. & 25 \\
\end{longtblr}
